\documentclass[10pt,a4paper]{amsart}
\usepackage[margin=19mm]{geometry}
\usepackage{amsmath,amssymb,mathtools}
\usepackage[scr=rsfs,cal=euler]{mathalfa}
\usepackage{xfrac}
\usepackage[unicode,hidelinks,bookmarks=false]{hyperref}
\newcommand{\ie}{i.e.\ }
\newcommand{\cf}{cf.\ }
\newcommand{\resp}{respectively\ }
\newcommand{\Subsection}{\subsection}
\providecommand{\nobreakdash}{\nobreak\hbox{-}}
\setlength{\emergencystretch}{2em}
\multlinegap0pt
\usepackage{amssymb}
\newtheorem{lemma}{Lemma}
\newtheorem{remark}{Remark}
\newtheorem{proposition}{Proposition}
\title{Compact Embedding Theorem Associated with Classical Weight Functions in Two Variables}
\author{M. K. NANGHO, B. J. NKWAMOUO and J.L. WOUKENG}
\date{Source: arXiv:2605.14732v3 (CC BY 4.0)}
\begin{document}
\maketitle

\noindent\emph{Source and licence.} Compact Embedding Theorem Associated with Classical Weight Functions in Two Variables. Version arXiv:2605.14732v3 (CC BY 4.0), distributed by arXiv under the Creative Commons Attribution 4.0 International licence, http://creativecommons.org/licenses/by/4.0/, as recorded in the arXiv licence metadata for this exact version and shown on the abstract page https://arxiv.org/abs/2605.14732v3. Full licence notice: \texttt{licenses/arxiv-2605.14732-CC-BY-4.0.txt}.\par\smallskip

\noindent\textit{Editorial scope.} Counted unit: the article's proposition dderiv (source0.tex lines 279--284). The article's complete original proof of it is delivered with it (source0.tex lines 285--311, one \texttt{proof} environment of 27 lines). No mathematics is written, repaired or re-derived: the statement and the proof are the article's own lines, in the article's own notation, and no retained line is substituted or re-flowed.  Retained uncounted as the source-local context the proof needs: lines 195--197; lines 254--259; lines 270--275.  Nothing was omitted from any retained span, so the package carries no omission record.  No retained line contains a citation, so the wrapper carries no bibliography and no bibliography entry.  No retained line contains a figure environment, an image-inclusion command or drawing code, so no image is delivered and the package declares no asset dependency.  Editorial formatting only, in addition: an AMS \texttt{amsart} preamble that restates the article's own declarations and macros used by the retained lines, namely the environments \texttt{lemma}, \texttt{remark}, \texttt{proposition} on separate counters, together with the standard packages those lines need.  The retained environments print in this document as Lemma 1, Remark 1, Proposition 1 in order of appearance, whereas the article prints the same items with the numbering of its own full document.  The complete native source of the article as distributed in the arXiv e-print for this exact version is bundled unchanged as \texttt{sources/200607/source0.tex}.
			\begin{equation}\label{bc}
				\lim\limits_{j\rightarrow\infty}1_{\partial\Omega_j}\left(\rho \Phi\,\nabla\,u\right)\cdot\overrightarrow{n_j}=0.
			\end{equation}

	{\begin{lemma}\label{inlem} \hspace{1cm}\\
			\begin{enumerate}
				\item Let  $u$ and $v$ be elements of $L^2(\Omega,\,\rho\Phi)$. Then, $\|u^t\Phi v\|_{L^1(\Omega,\,\rho)}\leq \|u\|_{L^2(\Omega,\,\rho\Phi)}\|v\|_{L^2(\Omega,\,\rho\Phi)}.$ 
				\item Let $u$ be an element of $L^2(\Omega,\,\rho)$. Then, for all $v\in\mathcal{P}\times\mathcal{P}$, there exists a constant $C\geq 0$ such that  $$\int_{\Omega}|u div(\rho\Phi )v|dx_1dx_2\leq C \| u\|_{L^2(\Omega,\,\rho)}.$$
			\end{enumerate}
		\end{lemma}

	\begin{remark}\label{bcremark}\hspace{1 cm}\\
		\begin{enumerate}
			\item If $\Omega$ is bounded, (\ref{bc}) becomes $1_{\partial\Omega}\left(\rho\Phi\nabla v\right)\cdot\overrightarrow{n}=0$, $v\in\mathcal{P}$.
			\item (\ref{bc}) can also be rewritten as 	$\lim\limits_{j\rightarrow\infty}1_{\partial\Omega_j}\left(\rho \Phi\,v\right)\cdot\overrightarrow{n_j}=0,\, v\in\mathcal{P}\times\mathcal{P}$.
		\end{enumerate}
	\end{remark}

	\begin{proposition}\label{dderiv}
		Let $u$ be a smooth function on $\overline{\Omega}$ such that $u\in L^2(\Omega,\,  \rho)$ and $\nabla u\in L^2(\Omega,\,  \rho\Phi)$. Then
		\begin{equation}\label{dderiv0}
			\int_{\Omega}u div(\rho\Phi v)dx_1dx_2=-\int_{\Omega}(\nabla u)^t\rho\Phi vdx_1dx_2,\,v=(v_1,v_2)^t\in \mathcal{P}\times\mathcal{P}.
		\end{equation}
	\end{proposition}

	\begin{proof}
		Observing that: \begin{equation}\label{deq}
			div(u\rho\Phi\,v)=udiv(\rho\Phi\,v)+\left(\nabla u\right)^t\rho\Phi\,v
		\end{equation} and integrating both sides on $\Omega$ we obtain 
		\begin{equation}\label{ideq}
			\int_{\Omega}udiv(\rho\Phi\,v)dx_1dx_2=\int_{\Omega}div(u\rho\Phi\,v)dx_1dx_2-\int_{\Omega}\left(\nabla u\right)^t\Phi\,v\rho dx_1dx_2.
		\end{equation}
		\begin{itemize}
		\item 	If $\Omega$ is bounded, from the divergence theorem, $\int_{\Omega}div(u\rho\Phi\,v)dx_1dx_2= {\int_{\partial\Omega}\left(\rho\Phi\nabla v\right)\cdot\overrightarrow{n}\,dx_1dx_2 }$. Using the boundary condition for bounded domain (cf. the first item of the Remark\, \ref{bcremark})  and the fact that  $u$ is defined on $\overline{\Omega}$, we have $1_{\partial\Omega}u\left(\rho\Phi v\right)\cdot\overrightarrow{n}=0$. Therefore, $\int_{\Omega}div(u\rho\Phi\,v)dx_1dx_2=0$.
		
		\item {	If $\Omega$ is unbounded, then \\$\displaystyle{\lim_{j \to \infty }1_{\Omega_j}div(u\rho\Phi\,v)}=1_{\Omega}div(u\rho\Phi\,v)$ and $\lvert1_{\Omega_j}div(u\rho\Phi\,v)\rvert\leq \lvert 1_{\Omega}div(u\rho\Phi\,v)\rvert$ because $\Omega_j$ is an {increasing} sequence of subsets of $\Omega$ such that $\bigcup\Omega_j=\Omega$.\\
		Combining the relation (\ref{deq}) and the Lemma \ref{inlem}, we derive the existence of a constant $C\geq 0$ such that
			$$\int_{\Omega}\lvert div(u\rho\Phi\,v)\rvert dx_1dx_2 \leq C \| u\|_{L^2(\Omega,\,\rho)}+\| \nabla u\|_{L^2(\Omega,\,\rho\Phi)}\| v\|_{L^2(\Omega,\,\rho\Phi)}.$$ Therefore, $\int_{\Omega}\lvert div(u\rho\Phi\,v)\rvert dx_1dx_2<+\infty$ for $u\in L^2(\Omega,\,\rho)$ and $\nabla u\in L^2(\Omega,\,\rho\Phi)$. So, by the dominated convergence theorem on the sequence $\left\{1_{\Omega_j}div(u\rho\Phi\,v)\right\}_{j\geq 1}$ it follows that:} 
		
		$$\int_{\Omega}div(u\rho\Phi\,v)dx_1dx_2=\lim_{j \to \infty }\int_{\Omega_j}div(u\rho\Phi\,v)dx_1dx_2.$$ Since  $\Omega_j,\,j=1,2,.. .$ is bounded, we deduce from the divergence theorem $$\int_{\Omega_j}div(u\rho\Phi\,v)dx_1dx_2=\int_{\partial\Omega_j}\left(u\rho\Phi\,v\right)\cdot\overrightarrow{n_j}dx_1dx_2.$$ Therefore, 
		
		$$\int_{\Omega}div(u\rho\Phi\,v)dx_1dx_2=\lim_{j \to \infty }\int_{{\Omega}} 1_{\partial\Omega_j}\left(u\rho\Phi\,v\right)\cdot\overrightarrow{n_j}dx_1dx_2.$$
		
		 From the boundary condition for unbounded domain (cf. the second item of the Remark\,\ref{bcremark}) $$\displaystyle{\lim_{j \to \infty}1_{\partial\Omega_j}\left(u\rho\Phi v\right)\cdot\overrightarrow{n_j}}=0$$ whenever $v\in \mathcal{P}\times \mathcal{P}$. So, to obtain (\ref{dderiv0}) it suffices to prove that  
		\begin{equation}\label{e10}
			\lim_{j \to \infty }\int_{{\Omega}} 1_{\partial\Omega_j}\left(u\rho\Phi\,v\right)\cdot\overrightarrow{n_j}dx_1dx_2=\int_{{\Omega}}  \lim_{j \to \infty } 1_{\partial\Omega_j}\left(u\rho\Phi\,v\right)\cdot\overrightarrow{n_j}dx_1dx_2.
		\end{equation} 
		Let  $\{f_j\}_{j\geq 1}$ {be the sequence of functions defined by} 
		$f_j=1_{\partial\Omega_j}\left(u\rho\Phi\,v\right)\cdot\overrightarrow{n_j},j=1,2,...\, .$ According to the Cauchy-Schwarz inequality, we deduce that:  $$|f_j|\leq 1_{\partial\Omega_j} |u|\|\overrightarrow{n_j}\|\|\Phi\,v\|\rho\leq 1_{\overline{\Omega}}|u|\|\Phi\,v\|\rho$$ for $\|\overrightarrow{n_j}\|=1$. Therefore, the use  of the {H\"older} inequality yields   
		$$\int_{\Omega}|u|\|\Phi\,v\|\rho dx_1dx_2\leq \left(\int_{\Omega}u^2\rho dx_1dx_2\right)^{1\over 2}\left(\int_{\Omega}\|\Phi\,v\|^2\rho dx_1dx_2\right)^{1\over 2}.$$ Since $u\in L^2(\Omega,\, \rho)$ and $\|\Phi\,v\|^2$ is a polynomial, the right-hand side is finite. So $\left\{f_j\right\}_{j\geq 1}$ is a sequence of integrable functions on $\Omega$. Moreover, there exists a positive integrable function, $h=|u|\|\rho\Phi\,v\|$, on $\Omega$ such that $|f_j|\leq h$. Therefore, from the dominated convergence theorem, (\ref{e10}) is satisfied. Thus $\int_{\Omega}div(u\rho\Phi\,v)dx_1dx_2=0.$ Using (\ref{ideq}), we obtain $$\int_{\Omega}udiv\left(\rho\Phi\right)dx_1dx_2=-\int_{\Omega}\left(\nabla u\right)^t\rho\Phi vdx_1dx_2\; v\in \mathcal{P}\times \mathcal{P}.$$
	\end{itemize}
	\end{proof}

\end{document}
